% ================= INTRODUCTION =================
\section{Introduction}

Classical blind motion deblurring recovers a sharp image without knowing the blur kernel by
reading the kernel off the blurred image's own spectrum: the angle of a linear motion blur from a
Radon transform of the log-spectrum and its length from the cepstrum, followed by non-blind
deconvolution such as Richardson--Lucy~\cite{richardson1972,lucy1974,cannon1976,moghaddam2007,oliveira2007}.
The approach needs no training data and every intermediate quantity can be inspected, but it was
developed and is nearly always evaluated on texture-rich natural photographs.

License plates and printed text are a different kind of input: small, low-contrast, dominated by a
few straight strokes and, in practice, resampled and compressed. Our proposal hypothesised that the
classical estimator is weakest exactly there, planned a text-contrast pre-processing stage to
strengthen the blur signature, and a confidence gate so that an untrustworthy estimate is refused
instead of silently making the image worse.

This report describes the full system and an evaluation that follows the proposal's five
objectives: (O1) the core pipeline validated against known kernels on synthetic blur and on the
K\"ohler camera-shake benchmark~\cite{kohler2012}; (O2) a controlled angle $\times$ length $\times$
noise sweep quantifying where it fails on natural images (Kodak) versus text and plates;
(O3) the pre-processing stage measured on the same sweep; (O4) the gate's trade-off between
harm avoided and good restorations withheld; and (O5) the full pipeline on real blurred plates from
public datasets (LPBlur~\cite{lpblur2024} and CCPD-Blur~\cite{ccpd2018}), with OCR accuracy before
and after restoration against the classical baseline.

\paragraph{Summary of findings.}
\begin{enumerate}
  \item \textbf{The difficulty is region size, noise and short blur, not text content.}
        In a plate-sized window printed text is the estimator's easiest input (98.9\% success),
        plates reach 89.4\% and natural texture 80.7\%; the same natural images reach 93.0\% at full
        frame. The gap the proposal expected is a \emph{small-region} gap --- exactly where plates
        live --- concentrated on short blurs under noise.
  \item \textbf{Edge enhancement does not strengthen the signature; band limiting does.} The planned
        edge-enhanced, locally normalised input halves estimation success (46.5\% vs.\ 89.7\%), and
        derivative operators in general collapse under noise; only a mild Gaussian high-pass helps on
        synthetic data (+2.5 points). On real plates the decisive operation is the opposite one:
        discarding the upper spectrum, where resampling and compression residue hide the nulls. The
        final text-aware stage is a multi-band search, which raises OCR accuracy on held-out real
        motion-blurred plates from 0.18 (blurred) and 0.24 (classical baseline) to 0.50.
  \item \textbf{The gate makes the pipeline safe to run.} Requiring a deep dip and agreement between
        image halves, it refuses 96--100\% of harmful restorations on held-out synthetic data; on
        LPBlur only 0.8\% of plates get worse. On CCPD-Blur, whose blur is mostly not linear motion,
        unconditional deblurring cuts accuracy from 0.76 to 0.19; the gate refuses 82\% of the
        images and recovers most of the loss (0.68), though it does not beat leaving them alone.
  \item \textbf{A correction to phase 0.} The Radon-variance angle is a more robust cue than an
        unrestricted joint (direction $\times$ lag) search; the final estimator uses the Radon angle
        with a $\pm3^\circ$ joint refinement, i.e.\ the proposal's pipeline as written.
\end{enumerate}
All numbers in this report are regenerated by \texttt{experiments/run\_all.py}; design choices were
made on development splits disjoint from every reported test image.
